% 相关工作
\section{相关工作}
\label{sec:related}

我们按 APT 检测逻辑组织相关工作，而非按具体模型架构分类。现有研究大体可以分为两类：一类从异常行为出发，学习正常系统行为并寻找偏离；另一类从攻击过程出发，试图恢复攻击路径或行为序列。本文工作与两类方法都相关：我们借鉴异常检测方法发现可疑起点，但把最终检测对象提升为因果链，并用链级模型判断该过程是否异常。

\paragraph{异常行为驱动的 APT 检测。}

第一类工作将 APT 检测视为异常发现问题：先学习良性系统行为的表示，再在图、时间窗口、节点、实体或边的粒度上寻找偏离。UNICORN~\cite{unicorn2020} 使用 Weisfeiler--Lehman 图核和直方图草图对整个溯源图进行分类；PROGRAPHER~\cite{yang2023prographer} 将溯源图划分为时序快照，用 graph2vec 和 TextRCNN 进行序列级异常检测。MAGIC~\cite{jia2024magic} 通过图掩码自编码器学习良性实体表示，并用 KNN 密度估计发现异常实体；ThreaTrace~\cite{wang2022threatrace} 和 FLASH~\cite{flash2024} 分别使用 GraphSAGE 或语义/图上下文编码进行节点级检测。KAIROS~\cite{cheng2024kairos} 进一步在边级计算时序图网络重建损失，并在时间窗口内聚合异常信号。EdgeTrace~\cite{edgetrace2025} 也以边级异常评分作为后续追踪的基础。

这类方法的优势是能够在大规模审计数据中有效定位“哪里不正常”：异常时间段、可疑实体、可疑节点或可疑边。但其检测对象通常仍是局部或聚合单元，而不是攻击调查真正需要的完整行为过程。因此，分析员往往还需要在告警之后手工判断这些局部异常之间是否存在因果联系。本文继承这一路线的异常检测能力：在当前实现中，我们使用 KAIROS 的边级评分作为种子来源；但不同之处在于，边级异常只用于提示“从哪里开始查”，而最终检测和评估发生在由因果约束提取出的事件链上。

\paragraph{攻击过程驱动的检测与调查。}

第二类工作直接面向攻击过程，试图把分散事件组织成路径、序列或子图。Holmes~\cite{holmes2019}、Poirot~\cite{poirot2019} 和 SLEUTH 通过专家定义的攻击模式在溯源图中追踪因果路径；SPARSE~\cite{sparse2024} 结合语义标签和威胁情报构造可疑语义图。近期学习方法也开始关注路径或序列层面的检测：SLOT~\cite{slot2025} 在异常检测后用标签传播聚类构造攻击链；CAGE~\cite{cage2025} 通过 Q-learning 为溯源边赋予动态因果权重，并提取加权攻击子图；EdgeTrace~\cite{edgetrace2025} 在边级异常评分之后进行语义聚类和路径推断。另一条相关路线是把溯源数据组织成序列后交给 Transformer 或类似架构建模，例如 EagleEye~\cite{eagleeye2024} 的固定窗口事件序列、Sentient~\cite{sentient2026} 的图 Transformer 与意图分析模块、GET-AID~\cite{getaid2025} 的时序子图注意力，以及 PanThreat~\cite{panthreat2025} 和 LogShield~\cite{logshield2024} 的日志/溯源序列建模。

这类方法更接近攻击调查需求，因为它们试图回答“攻击如何展开”。但其关键差异在于链或序列如何产生：有的方法依赖人工攻击模板或威胁情报，难以覆盖未知 APT；有的方法把路径重建放在节点/边级告警之后，链质量受上游漏报限制；还有的方法用时间窗口、随机游走或时序子图定义序列，未必保证相邻事件之间具有真实依赖关系。换言之，这些方法已经认识到过程级建模的重要性，但过程本身往往由模板、分割规则或后处理决定，而不是由一个可计算的因果链质量准则在检测前显式定义。

\paragraph{与本文工作的关系。}

本文属于攻击过程驱动的检测与调查方向，目标是把 APT 告警从局部异常推进到可解释的行为过程。与仅输出异常窗口、实体或边的方法不同，本文将最终检测对象定义为因果事件链；与依赖人工模板或告警后路径重建的方法不同，本文在检测之前提取结构可信的因果链，并在链级判断其是否偏离良性行为。本文的主要亮点在于把“因果链提取”和“链级异常检测”明确分层：前者保证检测对象具有可解释的依赖结构，后者利用自监督 Transformer 学习正常链级行为并识别语义异常。具体方法将在 \S\ref{sec:method} 中介绍。
